\section{Harmonic-tail reconstruction and complete even-order prefixes}
\label{opf:section}

Gamma residue calculations naturally produce products of Hurwitz tails.
Their centered even moments offer a second route to exact identities:
first replace the moments by Bernoulli coordinates, then reconstruct the
coordinates from a finite initial prefix. The determinant experiments in
the preserved reports \cite{opf:reports} lead to a uniform proof for each
fixed smaller order. We present the common reduction once. Complete
certificates close smaller orders ten and twelve from those reports,
and the independent extension below closes smaller order fourteen.

\subsection{Centered continuation and the coefficient model}

Fix integers $1\le a<b$, put $p=2a$, $q=2b$, and define
\[
 A_s(n)=\zeta(s,n+1)=\zeta(s)-H_n^{(s)},\qquad x_n=n+\tfrac12,
 \qquad \mathcal D_{p,q}(w)=\sum_{n\ge0}A_p(n)A_q(n)x_n^{-w}.
\]
The defining half-plane is $\Re w>3-p-q$. The centered expansion
\begin{equation}\label{opf:tail-expansion}
 A_s(n)\sim\frac{x_n^{1-s}}{s-1}
 +\sum_{h\ge1}\frac{B_{2h}(1/2)}{(2h)!}(s)_{2h-1}
                                x_n^{1-s-2h}
\end{equation}
continues $\mathcal D$ by finitely many subtractions on every compact
parameter set. Its possible poles are odd integers
$3-p-q-2h$, so $D_r=\mathcal D_{p,q}(-2r)$ is regular for every
integer $r\ge0$. Equivalently, $D_r$ is the constant coefficient in
the full large-$N$ expansion of the partial sum through $n=N-1$.
Indeed, every nondecaying subtraction is an even power of $x_n$,
whose partial sum is an odd polynomial in $N$; its continued sum
is $\zeta(-2h,1/2)=0$. This fixes the cutoff coordinate.

Define the triangular transforms
\begin{align}
 C_j&=\sum_{r=0}^j\binom{2j+1}{2r}2^{-2(j-r)}D_r,
 \label{opf:C}\\
 D_r&=\sum_{j=0}^r\frac{\binom{2r}{2j}}{2j+1}
                                  B_{2r-2j}(1/2)C_j.
 \label{opf:D}
\end{align}
The Bernoulli generating function proves the inverse formula
coefficientwise. Let $Z_k=\zeta(2k)$, $0\le k<b$, with $Z_0=-1/2$.
For every $j\ge b$ the exact late-coordinate identity is
\begin{equation}\label{opf:late}
 C_j=\sum_{k=0}^{b-1}M_{j-b,k}Z_k,
\end{equation}
where, for all $t\ge0$ and $0\le k<b$,
\begin{equation}\label{opf:M}
 M_{tk}=B_{2(t+k-a+1)}
 \left\{\frac{\binom{2t+2}{2a-2k}}{2t+2}
 +\frac{\binom{2b+2t-2a+2}{2b-2k}}{2b+2t-2a+2}\right\}.
\end{equation}
Negative Bernoulli indices give zero. Each binomial summand is
omitted unless its lower index lies between two and its upper index.
In particular, a lower index zero is omitted. Write $M_{a,b}$ for
the square block $0\le t,k<b$.

Here is a direct analytic justification of \eqref{opf:late}. Put
$h=2j+1$. Summation by parts at the finite cutoff gives
\begin{align*}
 C_j=\operatorname{CT}_{N}\biggl[&N^h A_p(N)A_q(N)
 +\sum_{m=1}^N m^{h-q}A_p(m)
 +\sum_{m=1}^N m^{h-p}A_q(m)
 +\sum_{m=1}^N m^{h-p-q}\biggr].
\end{align*}
The two positive odd exponents $h-q,h-p$ permit Faulhaber expansion.
For an odd positive integer $r$,
\[
 \sum_{m=1}^{n-1}m^r=
 \frac1{r+1}\sum_{\ell=0}^{r}\binom{r+1}{\ell}B_\ell
                                                  n^{r+1-\ell}.
\]
There is no additional constant term: the $B_{r+1}$ in the
Bernoulli-polynomial formula cancels. The two $B_1=-1/2$ terms
cancel the final harmonic sum. The even Bernoulli terms give
exactly \eqref{opf:M}, including the omitted lower-index-zero terms.
More explicitly, after that cancellation the endpoint remainder is
\[
 \begin{aligned}
 &N^h A_p(N)A_q(N)+Q_{h-q}(N)A_p(N)+Q_{h-p}(N)A_q(N)\\
 &\quad-\sum_{(s,r)=(p,h-q),(q,h-p)}\ \sum_{\substack{0\le\ell\le r\\\ell\text{ even}}}
 \frac{\binom{r+1}{\ell}B_\ell}{r+1}A_{s-r-1+\ell}(N).
 \end{aligned}
\]
Here tails with nonpositive even order mean their Hurwitz-polynomial
continuations. This remainder has zero constant coefficient.
To check this without discarding a finite counterterm, use its
formal Euler--Maclaurin series: for even $s$,
$\mathsf A_s(-N)=-\mathsf A_s(N)-N^{-s}$, and for odd $r$,
$Q_r(-N)=Q_r(N)-N^r$, where $Q_r(N)=\sum_{m=1}^N m^r$.
Substitution makes the combined remainder odd in $N$. These are
identities of asymptotic Laurent series, not evaluations of the
Hurwitz tail at a negative integer. A finite number of their terms
suffices for every fixed $j$, by the remainder bound in
\eqref{opf:tail-expansion}. Thus no divergent interchange is used.

Before evaluating these coordinates, regard
$C_0,\ldots,C_{b-1},Z_0,\ldots,Z_{b-1}$ as a formal rational
coordinate list. The first $2b=q$ moment rows have determinant
\begin{equation}\label{opf:prefix-determinant}
 \det M_{a,b}/(4b-1)!!.
\end{equation}
This follows from the block diagonal matrix $\operatorname{diag}(I_b,M_{a,b})$
for the $C$ rows and the diagonal entries $1/(2r+1)$ of
\eqref{opf:D}. Formal coefficient rank is the claim being proved;
it does not imply arithmetic independence of the evaluated constants.
Every rational coefficient identity nevertheless remains valid after
evaluation.

\begin{theorem}[Effective eventual reconstruction]\label{opf:main}
For every fixed integer $a\ge1$, a terminating exact algorithm
constructs an integer $B_a$ for which $\det M_{a,b}\ne0$ at every
integer $b\ge B_a$. Consequently the complete fixed-order prefix
question reduces to the finite set $a<b<B_a$.
\end{theorem}

The following normalization and threshold proof are the common
argument of the quartic reconstruction report \cite{opf:reports}.
We retain the full scalar multiplier and finite Pad\'e calculation,
which make the later certificates certificates of the original
matrix rather than of an experimentally related surrogate.

\subsection{Exact normalization of the retained determinant}

Assume initially $b\ge2a$. For $0\le i,\ell<a$ define
\[
 n_i=b-a+i,\quad N_{i\ell}=2(n_i-\ell)+1,\quad
 R_\ell=2b-2a+2+2\ell,
\]
\[
 F_{i\ell}=\frac1{(2\ell+1)!}
 \prod_{r=0}^{2\ell-1}(2n_i-r)(2n_i+3+r),
\]
and define integer-polynomial coefficients $d_{ih}$ by
\begin{equation}\label{opf:newton}
 (2n_i+2-y)(4n_i+2-y)^{\underline{2i}}
       =\sum_{h=0}^{2i+1}d_{ih}y^{\underline h}.
\end{equation}
They lie in $\mathbb Z[b]$: the change of basis from ordinary
powers to falling powers has integer Stirling-number coefficients.
Let $U_1,\ldots,U_{2a-1}$ be independent variables and put
\begin{equation}\label{opf:A}
 \mathcal A_{i\ell}=2F_{i\ell}+
 \frac{R_\ell!}
 {(2n_i+1)(2n_i+2)(2\ell+2)!N_{i\ell}!}
 \sum_{k=0}^{i}d_{i,2k+1}N_{i\ell}^{\underline{2k+1}}
            U_{a-i+\ell+k}.
\end{equation}
The inherited reduction is the theorem ``Reduction to a determinant
of size $a$,'' equation \texttt{eq:general-small-det}, in the source
file \texttt{03b\_general\_hankel.tex}~\cite{dpm:reports,opf:reports}.
We give its exact multiplier and the normalization argument, so
that no nonvanishing statement about an uncomputed inverse matrix
is needed.

\begin{lemma}[Exact source-to-retained determinant multiplier]
\label{opf:normalization}
Put
\[
 h_n=\frac{((2n+1)!)^2}{(4n+2)!(4n+3)!},\qquad
 c_n=\frac{(2n+1)(2n+2)!}{(4n+2)!},
 \qquad \epsilon_{a,b}=a(b-a)+\frac{a(a-1)}2.
\]
For $b\ge2a$, with $U_j=2B_{2b-2j}$,
\begin{equation}\label{opf:exact-multiplier}
 \det M_{a,b}=(-1)^{\epsilon_{a,b}}
 \frac{\prod_{t=0}^{b-1}(2b+2t-2a+1)!}
      {\prod_{k=0}^{b-1}(2b-2k)!}
 \left(\prod_{n=0}^{b-a-1}h_n\right)
 \left(\prod_{i=0}^{a-1}c_{n_i}\right)\det\mathcal A.
\end{equation}
Every scalar factor other than the displayed sign is strictly
positive. In particular,
\begin{equation}\label{opf:equivalence}
 \det M_{a,b}\ne0
 \quad\Longleftrightarrow\quad
 \det\mathcal A\bigl(b;2B_{2b-2},\ldots,2B_{2b-4a+2}\bigr)\ne0.
\end{equation}
\end{lemma}

\begin{proof}
Set $\mu_j=B_{2j}/(2j)!$ for $j\ge0$ and $\mu_j=0$ for $j<0$.
Cancelling factorials in \eqref{opf:M} gives
\[
 M_{a,b}=R(H+U^{\mathrm{pert}})E,
 \quad R_{tt}=(2b+2t-2a+1)!,\quad E_{kk}=\frac1{(2b-2k)!},
\]
where $H_{tk}=\mu_{t+k-a+1}$ and
\[
 U^{\mathrm{pert}}_{tk}=\mathbf1_{k<a}\mu_{t+k-a+1}
 \frac{(2t+1)!(2b-2k)!}
 {(2a-2k)!(2b+2t-2a+1)!}.
\]
The diagonal symbols $R,E$ in this proof are matrices; $R_\ell$
in \eqref{opf:A} is the previously defined integer.

The explicit monic Pad\'e rows are
\[
 \pi_n(x)=\sum_{t=0}^np_{nt}x^t,\qquad
 p_{nt}=\frac{2(2n+2t+1)!(2n+1)!}
 {(4n+2)!(2n-2t+1)!(2t)!}.
\]
For $\mathcal N(x^r)=\mu_{r+1}$ they satisfy
\begin{equation}\label{opf:pade-orthogonality}
 \mathcal N(x^j\pi_n)=0\ (0\le j<n),\qquad
 \mathcal N(x^n\pi_n)=h_n.
\end{equation}
Here is a finite verification of the normalization. Define
\[
 P_N(t)=\sum_{k=0}^N
 \frac{(2N-k)!N!}{(2N)!k!(N-k)!}t^k.
\]
Direct coefficient comparison gives
\[
 e^tP_N(-t)-P_N(t)
 =\frac{(-1)^N(N!)^2}{(2N)!(2N+1)!}t^{2N+1}
       +O(t^{2N+2}).
\]
Indeed, the coefficients through $2N$ reduce to the finite identity
$\sum_k(-1)^k\binom Nk(1+x)^{2N-k}=x^N(1+x)^N$;
the next coefficient is
$(2N)!^{-1}\int_0^1x^N(x-1)^N\,dx$.
Taking $N=2n+1$, set
\[
 q_n(t^2)=\frac{P_N(t)-P_N(-t)}t,\qquad
 r_n(t^2)=\frac{P_N(t)+P_N(-t)}2.
\]
For $\Phi(z)=\sum_{j\ge0}\mu_jz^j
=\frac{\sqrt z}{2}\coth\frac{\sqrt z}{2}$ this proves
\[
 \Phi(z)q_n(z)-r_n(z)=h_nz^{2n+1}+O(z^{2n+2}),
 \qquad \pi_n(x)=x^nq_n(1/x).
\]
Coefficient comparison proves \eqref{opf:pade-orthogonality};
it also gives, for $n\ge\ell\ge0$,
\begin{equation}\label{opf:constant-pade-row}
 \sum_{t=0}^np_{nt}\mu_{t-\ell}=c_n
 \frac1{(2\ell+1)!}
 \prod_{r=0}^{2\ell-1}(2n-r)(2n+3+r).
\end{equation}
Specifically, this sum is $[z^{n-\ell}]r_n(z)$, and its
factorial coefficient formula is the right side. Thus both
$c_n$ and $h_n$ have been obtained without an assumed Pad\'e
normality assertion.

Replace row $n$ of $H+U^{\mathrm{pert}}$ by the linear combination
with coefficients $p_{nt}$, for $0\le n<b$. This simultaneous
transformation is unit lower triangular. In columns $k\ge a$ the
result is $\mathcal N(x^{k-a}\pi_n)$. Hence the final $a$ rows
vanish there, while the first $b-a$ rows in those columns form an
upper triangular square matrix with diagonal
$h_0,\ldots,h_{b-a-1}$. The retained $a$-square matrix is
\[
 A^{\mathrm{ret}}_{ik}=
 \sum_{t=0}^{n_i}p_{n_it}\mu_{t+k-a+1}
 \left(1+\frac{(2t+1)!(2b-2k)!}
 {(2a-2k)!(2b+2t-2a+1)!}\right).
\]
Moving the first $a$ columns past the others contributes
$(-1)^{a(b-a)}$. Reverse the retained columns, putting
$\ell=a-1-k$, and divide row $i$ by $c_{n_i}$.
This adds the determinant factor
$(-1)^{a(a-1)/2}\prod_i c_{n_i}$.

It remains to identify this divided matrix with \eqref{opf:A}.
Its unperturbed entry is $F_{i\ell}$ by
\eqref{opf:constant-pade-row}. Writing $n=n_i$, $N=N_{i\ell}$,
and
\[
 W(j)=(j+2\ell+1)(2n+2\ell+j+1)^{\underline{2i}},\qquad
 C=\frac{R_\ell!}{(2n+1)(2n+2)(2\ell+2)!N!},
\]
factorial cancellation gives the perturbation entry
$2C\sum_{r=0}^{n-\ell}\binom N{2r}B_{2r}W(2r)$.
Since $W(N-y)$ is exactly the polynomial in \eqref{opf:newton},
use the finite identity
\[
 \sum_{j=0}^N\binom Nj B_j(N-j)^{\underline h}
    =N^{\underline h}B_{N-h}(1).
\]
It follows immediately by taking $N^{\underline h}$ outside
the binomial sum and using the defining translation formula for
Bernoulli polynomials. Here $N\ge2i+3$, so all indices $N-h$
are at least two. Only odd $h$ survive. Restoring the omitted
$j=1$ term, with $B_1=-1/2$, gives
\[
 \sum_r\binom N{2r}B_{2r}W(2r)
  =\sum_{k=0}^i d_{i,2k+1}N^{\underline{2k+1}}
           B_{2b-2(a-i+\ell+k)}+\frac N2W(1).
\]
Finally, $CNW(1)=F_{i\ell}$ by factorial cancellation. The divided
retained entry is therefore $2F_{i\ell}$ plus the Bernoulli term
in \eqref{opf:A}. Restoring all row and column factors proves
\eqref{opf:exact-multiplier}.
\end{proof}

For reference, the unperturbed base determinant has the exact
normalization
\begin{equation}\label{opf:base-determinant}
 \det H=(-1)^{\epsilon_{a,b}}
 \left(\prod_{n=0}^{b-a-1}h_n\right)
 \left(\prod_{i=0}^{a-1}c_{n_i}\right)\det F\ne0.
\end{equation}
One can verify the last inequality without using it in the lemma.
Conjugate $H$ by $\operatorname{diag}(1,-1,\ldots)$.
Euler's Bernoulli formula shows that entries with $t+k\ge a$
come from a strictly definite positive moment matrix,
with sign $(-1)^a$. For $a=2h$ the resulting block matrix is
$\left(\begin{smallmatrix}0&B\\B^{\mathsf T}&P\end{smallmatrix}\right)$,
where $P$ is positive definite and $B$ has full row rank $h$.
For $a=2h+1$ it is
$\left(\begin{smallmatrix}ee^{\mathsf T}&B\\B^{\mathsf T}&-P\end{smallmatrix}\right)$,
where $P$ is positive definite, $e=(0,\ldots,0,1)^{\mathsf T}$,
and the first $h$ rows of $B$ have a nonsingular antitriangular
minor. The Schur complements are respectively strictly negative
and strictly positive. The inertia is $(b-h,h)$ in the even case
and $(h+1,b-h-1)$ in the odd case. Together with
\eqref{opf:base-determinant}, this also proves $\det F>0$.
Strict positivity of the moment blocks follows from
$\sum_{m\ge1}(2\pi m)^{-2s}P_0((2\pi m)^{-2})^2>0$ for a
nonzero polynomial $P_0$ and $s\ge1$.

\subsection{The general highest-degree determinant identity}

All denominators can be cleared constructively. Set
\begin{align*}
 J_a(b)&=\prod_{i=0}^{a-1}(2n_i+1)(2n_i+2),\\
 E_a(b)&=\prod_{i,\ell=0}^{a-1}
           \prod_{r=1}^{\max\{0,2i-4\ell-1\}}(R_\ell+r),\\
 D_a(b)&=(2a)!J_a(b)E_a(b),\qquad
 e_a=\sum_{i,\ell=0}^{a-1}\max\{0,2i-4\ell-1\}.
\end{align*}
The symbols $E_a(b)$ and $D_a(b)$ here are auxiliary polynomials;
they are unrelated to the moment coordinates.
Because $N_{i\ell}-R_\ell=2i-4\ell-1$, the polynomial $D_a$
clears every negative factorial ratio in \eqref{opf:A}, every
row denominator, and every constant factorial denominator.
Consequently
\[
 \widetilde{\mathcal A}=D_a\mathcal A
       \in\mathbb Z[b,U_1,\ldots,U_{2a-1}]^{a\times a},
 \qquad \mathsf P_a=\det\widetilde{\mathcal A}
       \in\mathbb Z[b,U_1,\ldots,U_{2a-1}].
\]
Also $D_a(b)>0$ for $b\ge2a$.

Put
\[
 F_0=(2b-4a+2)!,\qquad
 V_j=U_j\frac{F_0}{(2b-2j)!},\qquad
 \kappa_a=\frac{a!((2a)!)^a}{\prod_{\ell=0}^{a-1}(2\ell+2)!},
\]
and
\begin{align}
 \Gamma_a(b)={}&\kappa_aJ_a(b)^{a-1}E_a(b)^a
 \prod_{i=0}^{a-1}(4n_i+2)^{\underline{2i}}
 \prod_{\ell=0}^{a-1}\prod_{r=1}^{2a+2\ell}(2b-4a+2+r).
 \label{opf:Gamma}
\end{align}
The subscript $[a]$ denotes the part homogeneous of total degree
$a$ in the $U$ variables. The exact factorization is
\begin{equation}\label{opf:leading}
 (\mathsf P_a)_{[a]}
 =(-1)^{a(a+1)/2}\Gamma_a(b)
           \det(V_{i+j+1})_{0\le i,j<a}.
\end{equation}
To prove it, write $T=\mathcal A-2F$ and substitute the $V$ variables.
For the term $U_{a-i+\ell+k}$,
$N_{i\ell}-(2k+1)=2b-2(a-i+\ell+k)$, so all nonconstant
factorials cancel. Exactly as in the order-eight proof,
\[
 T=\operatorname{diag}\!\left(\frac1{(2n_i+1)(2n_i+2)}\right)
    L\,(V_{a-i+\ell})_{0\le i,\ell<a}
    \operatorname{diag}\!\left(\frac{R_\ell!}{F_0(2\ell+2)!}\right),
\]
where $L_{i,i-k}=d_{i,2k+1}$ and
$d_{i1}=-(i+1)(4n_i+2)^{\underline{2i}}$.
The $a$ diagonal signs of $L$, together with row reversal of the
middle matrix, give $(-1)^{a(a+1)/2}$. The remaining products are
precisely \eqref{opf:Gamma} after multiplication by $D_a^a$.
This proves \eqref{opf:leading} for every $a$, without evaluating
any large determinant experimentally.

At $U_j=2B_{2b-2j}$, Euler's Bernoulli formula gives
\begin{equation}\label{opf:actual-leading}
 \left|(\mathsf P_a)_{[a]}\right|
 =\frac{\Gamma_a(b)(4F_0)^a}{(2\pi)^{2a(b-a)}}
   \det\bigl(\zeta(2b-4a+2+2i+2j)\bigr)_{0\le i,j<a}.
\end{equation}
Its sign is $(-1)^{ab+a(a+1)/2}$. Define the positive rational
number
\[
 V_a=\prod_{1\le r<s\le a}(r^{-2}-s^{-2})^2,
 \qquad V_1=1.
\]
Keeping the first $a$ support points in Cauchy--Binet proves
\begin{equation}\label{opf:moment-lower}
 \det\bigl(\zeta(2b-4a+2+2i+2j)\bigr)
       \ge V_a(a!)^{-2b+4a-2}>0.
\end{equation}
Thus the highest-degree term is always nonzero in the stated range.

\subsection{An exact algorithm for a sufficient threshold}

Let $\|f\|_1$ denote the sum of the absolute values of all integer
coefficients of a polynomial in $b,U_1,\ldots,U_{2a-1}$. Compute
\[
 E_a^*=\max_{0\le i,\ell<a}\|
                  \widetilde{\mathcal A}_{i\ell}\|_1,
 \qquad K_a=a!(E_a^*)^a,
 \qquad g_a=6a^2-4a+ae_a.
\]
These are finite integer calculations on only $a^2$ explicitly given
entries; no expansion of the whole determinant is required.
The coefficient norm is submultiplicative, so the full coefficient
norm of $\mathsf P_a$ is at most $K_a$.

Each scaled entry has degree at most $6a+e_a-4$ in $b$.
For the nonconstant part this follows from
$\deg_b d_{ih}\le2i+1-h$, the degree $h$ of
$N_{i\ell}^{\underline h}$, and the degree difference
$R_\ell-N_{i\ell}=1+4\ell-2i$; the total bound is
$\deg D_a+4\ell$. The same bound holds for $D_aF_{i\ell}$.
Hence $\deg_b\mathsf P_a\le g_a$.
The factored polynomial \eqref{opf:Gamma} also has degree exactly
$g_a$. Every linear factor in it is at least $b$ when $b\ge4a$,
so
\[
 \Gamma_a(b)\ge\kappa_a b^{g_a}\quad(b\ge4a).
\]

For $b\ge4a+10$, define
\[
 T_{a,b}=\frac{8(2b)^{4a-4}F_0}{(2\pi)^{2b-4a+2}}\ge1.
\]
For the last inequality, put $s_0=2b-4a+2\ge26$;
$26!>8^{26}$ implies $s_0!>8^{s_0}$ by induction,
and $2\pi<8$.
As before, $|2B_{2b-2j}|\le T_{a,b}$ for $1\le j\le2a-1$.
Since the remaining terms have total $U$ degree at most $a-1$,
\[
 |(\mathsf P_a)_{[\le a-1]}|
      \le K_a b^{g_a}T_{a,b}^{a-1}.
\]
Combining this with \eqref{opf:actual-leading} and
\eqref{opf:moment-lower}, cancelling the common powers of $2\pi$,
and then using $1<2\pi<8$, gives
\begin{equation}\label{opf:bound}
 \frac{|(\mathsf P_a)_{[\le a-1]}|}{|(\mathsf P_a)_{[a]}|}
 \le R_a(b):=C_a\frac{b^{L_a}(8a!)^{2b}}{(2b-4a+2)!},
\end{equation}
where every constant is explicit and rational:
\begin{align}
 L_a&=4(a-1)^2,\qquad \chi_a=2a^2-6a+2,\notag\\
 C_a&=\frac{K_a\,8^{a-1+\max(\chi_a,0)}2^{L_a}}
 {\kappa_a4^aV_a(a!)^{4a-2}}.
 \label{opf:constant}
\end{align}
The use of $8^{\max(\chi_a,0)}$ also covers $a=1,2$, where
$\chi_a$ is negative.

Here is a terminating purely arithmetic prescription. Start at
\[
 B_a^{(0)}=\max\{4a+10,\ L_a,\ 16a!\},
\]
and increase the integer $b$ until the exact rational comparison
$R_a(b)<1$ holds. Define $B_a$ to be that first integer. The procedure
must terminate, because for every $b\ge B_a^{(0)}$,
\[
 \frac{R_a(b+1)}{R_a(b)}
 =\left(1+\frac1b\right)^{L_a}
 \frac{(8a!)^2}{(2b-4a+4)(2b-4a+3)}<\frac34.
\]
Indeed $b\ge L_a$ gives $(1+1/b)^{L_a}<3$,
$b\ge4a$ makes both denominator factors at least $b$, and
$b\ge16a!$ supplies the last numerical factor.
The comparison uses integer factorials and powers after clearing
the fixed rational denominator of $C_a$.

For every $b\ge B_a$, \eqref{opf:bound} is strictly less than one,
so the highest-degree term cannot be cancelled. Equations
\eqref{opf:equivalence} and $D_a(b)>0$ prove
Theorem~\ref{opf:main}. In particular, for each input $a$ one may
inspect the finite set $a<b<B_a$ exactly: this either proves the
entire fixed-order family or produces a genuine counterexample.
The construction does not perform those finite calculations for all
$a$, and does not settle the unrestricted conjecture.


\subsection{Complete smaller orders ten, twelve and fourteen}

\begin{theorem}[Three complete fixed-order families]\label{opf:complete}
For $a\in\{5,6,7\}$ and every integer $b>a$, $\det M_{a,b}\ne0$.
For every even $q>p$ with $p\in\{10,12,14\}$, the first $q$
centered moments $D_0,\ldots,D_{q-1}$ therefore have formal rank $q$.
Every later $D_r$ is an explicitly constructible rational combination
of that same prefix.
\end{theorem}

\begin{proof}
The eventual theorem gives sufficient thresholds
$B_5=1920$, $B_6=11520$, and $B_7=80640$.
For $a=7$, exact expansion of the $49$ scaled entries in
\eqref{opf:A} gives $e_7=56$, $g_7=658$, $L_7=144$, and
\[
 E_7^*<2^{315}.
\]
The stored exact rational comparison \eqref{opf:bound} is less than one
already at $B_7^{(0)}=16\cdot7!=80640$. The complete coefficient
construction, factorial cancellations and comparison are regenerated
by \path{verification/prefix-order-fourteen/eventual_threshold.py}.
For $a=6$ the boundary-product report also proves the simpler
entry bound $E_6^*<2^{512}$ by counting factors in the scaled entries.
Substituting that bound in \eqref{opf:constant} gives the same
sufficient $B_6=11520$; its exact integer comparison is retained
with the full replay. These calculations instantiate the written
infinite theorem; they do not infer it from samples.

At order ten, the alternating-tail certificate proves the uniform
congruence $2B_{2m}\equiv\mathcal B(m)\pmod{2^{16}}$ for $m\ge12$,
where
\[
 \mathcal B(m)=1+20068m-26384m^2+2816m^3-19712m^4
                                      +4096m^5-24576m^6.
\]
For completeness, this uniform congruence comes from Faulhaber's
identity for the sum of $2m$th powers below $256$. The $h=2$
correction vanishes, and the corrections $h\ge3$ have valuation
at least $8h-v_2(h)-1\ge23$, using the standard Bernoulli
denominator bound $v_2(B_{2r})\ge-1$.
After division by $128$, the even summands vanish when $m\ge12$.
For odd $v$, expand $v^{2m}=(1+(v^2-1))^m$; since $8\mid v^2-1$,
terms beyond degree seven vanish at the required precision.
Exact coefficient comparison of
$128^{-1}\sum_{v=1,3,\ldots,255}\sum_{j=0}^7
\binom{X}{j}(v^2-1)^j$ gives $\mathcal B(X)$ modulo $2^{16}$,
with rational coefficient valuations checked before any reduction.

Let $P_{10}(b;U)$ be the normalized integral numerator of
$\det\mathcal A/\det F$ fixed by the preserved 439-group
coefficient certificate. Its defining identity is
\[
 \frac{\det\mathcal A}{\det F}
 =\frac{P_{10}(b;U)}{138240(2b-7)(2b-5)D_{10}(b)},
\]
where the integer polynomial $D_{10}$ is specified by
\begin{align*}
 \det F={}&\frac{16384}{4465125}(b-3)(b-2)^2(b-1)
                  (2b-7)(2b-5)^2(2b-3)\\
 &\cdot(4b-15)(4b-13)(4b-11)^2(4b-9)^2\\
 &\cdot(4b-7)^2(4b-5)(4b-3)D_{10}(b).
\end{align*}
The finite coefficient verification clears row denominators, sums
all $120$ signed permutation products and checks this defining
identity coefficientwise. It also checks that $D_{10}(b)$ is odd
at every integer $b$, and the other displayed factors are nonzero
for $b\ge10$. Its finite polynomial numerator substitution is
$P_{10}(b;\mathcal B(b-1),\ldots,\mathcal B(b-9))
\equiv32R(b)\pmod{2^{16}}$, with
\begin{align*}
 R(b)={}&-512b^8+128b^7-928b^6-816b^5+464b^4\\
        &+352b^3+118b^2+217b+694.
\end{align*}
The finite determinant expansion preserves all $10403$ coefficients.
Every coefficient beyond the linear one is even and that linear
coefficient is odd, so $v_2(R(x)-R(y))=v_2(x-y)$. The unique root
modulo $2048$ is $962$. Thus, for $b\ge21$ outside that class,
the numerator is nonzero. Exact source matrices cover $6\le b\le20$.
The only remaining integer below $1920$ is $962$. The
boundary-product report computes
\[
 \det\mathcal A(5,962)\equiv60407\pmod{65537},\qquad
 \det M_{5,962}\equiv15578\pmod{65537}.
\]
All factorial and Bernoulli denominators are units at that prime.
Our independent rational reconstruction additionally checks the
five-square matrix, both inverse products, and its nonzero determinant;
a separate Wolfram calculation confirms every rational entry and
the determinant. The congruence proof, exceptional certificate and
eventual theorem cover every $b\ge6$. The two order-ten completions
in the newer reports retain credit for this synthesis.

At order twelve, the complete certificate gives nonzero residues
modulo the prime $65537$ for all $11513$ integers $7\le b\le11519$.
At order fourteen, our certificate gives nonzero residues modulo
the prime $7340033$ for all $80632$ integers $8\le b\le80639$.
In each case use the original $b$-square matrix for $b<2a$ and
\eqref{opf:A} otherwise. The prime exceeds every factorial index
and every cleared denominator occurring in this finite range;
Bernoulli numbers are computed as the coefficients of
$t/(e^t-1)$ over the field. A nonzero modular determinant therefore
proves the corresponding rational determinant is nonzero.
Together with \eqref{opf:exact-multiplier} and the eventual ranges,
this covers every integer $b>a$ in all three cases.

For reproducibility, the order-twelve full replay freshly builds
the native C++ verifier and compares its entire residue list with
the independent Python implementation. The order-fourteen replay
computes all $161281$ Bernoulli coefficients twice, using FLINT
series inversion and exact integer Kronecker convolution with a
complete product identity. Its native forward-difference/Gaussian
elimination determinants agree on every case with a separate
symbolic Stirling construction and division-free subset expansion.
All convolution and signed integer accumulation bounds are checked
before use; no floating-point rank or determinant enters the proof.
The complete residue lists and source programs are preserved.
Finally, \eqref{opf:late} and \eqref{opf:D} reconstruct every later
moment from the now invertible prefix.
\end{proof}

\subsection{Ordinary convergent identities from the reconstructed coordinates}

The rank result has an analytic payoff at every larger even order.
For any pair covered by Theorem~\ref{opf:complete} and $0\le k<b$, set
\begin{align}
 P_{a,b,k}(x)&=\sum_{t=0}^{b-1}(M_{a,b}^{-1})_{kt}
 \bigl\{(x+\tfrac12)^{2b+2t+1}-(x-\tfrac12)^{2b+2t+1}\bigr\}
 =\sum_{r=0}^{2b-1}w_{a,b,k,r}x^{2r},\label{opf:weight}\\
 s_{p,0}&=1/(p-1),\qquad
 s_{p,h}=B_{2h}(1/2)(p)_{2h-1}/(2h)!\quad(h\ge1),\notag\\
 \alpha_h&=\sum_{v=0}^h s_{p,v}s_{q,h-v},\notag\\
 S_{a,b,k}(x)&=\sum_{r=a+b-1}^{2b-1}w_{a,b,k,r}
           \sum_{h=0}^{r-a-b+1}\alpha_h x^{2(r-a-b+1-h)}.
 \label{opf:subtraction}
\end{align}
Then the following are ordinary absolutely convergent sums:
\begin{equation}\label{opf:ordinary}
 \boxed{\sum_{n\ge0}\left\{
 P_{a,b,k}(x_n)\frac{\psi^{(p-1)}(n+1)\psi^{(q-1)}(n+1)}
                         {(p-1)!(q-1)!}-S_{a,b,k}(x_n)
                    \right\}=\zeta(2k).}
\end{equation}
Indeed, \eqref{opf:subtraction} removes every nondecaying term
of the centered product expansion, leaving $O(x_n^{-2})$.
The polynomial difference in \eqref{opf:weight} is exactly the
$C$ transform, and the inverse matrix makes its continued sum $Z_k$.
Each subtracted even monomial has zero centered continued sum,
so the remaining ordinary sum has that same value. For even $p$,
$\psi^{(p-1)}(n+1)=(p-1)!A_p(n)$, which gives the displayed form.
This construction supplies exact identities with rational finite
counterterms for all the stated pairs, including smaller order fourteen.

The unrestricted prefix conjecture still requires every fixed
$a\ge8$ finite completion (and any other unaccepted smaller-order
branches). No arithmetic independence or Gaussian period identity
follows from the coefficient theorem. The proof combines a written
analytic reduction, an infinite nonvanishing argument, and complete
exact finite certificates, with distinct roles for each.
